\documentclass[10pt]{article}
\usepackage[T1]{fontenc}
\usepackage{lmodern,microtype,amsmath,amssymb,amsthm,mathtools,tikz-cd}
\usepackage[margin=1in]{geometry}
\usepackage{xurl,hyperref}
\hypersetup{colorlinks=true,linkcolor=blue,citecolor=blue,urlcolor=blue,
pdftitle={Internal theories and intrinsic incompleteness},pdfauthor={Florian Lengyel}}
\newtheorem{theorem}{Theorem}[section]
\newtheorem{proposition}[theorem]{Proposition}
\newtheorem{lemma}[theorem]{Lemma}
\newtheorem{corollary}[theorem]{Corollary}
\newcommand{\I}{\mathcal I}
\newcommand{\C}{\mathcal C}
\newcommand{\D}{\mathcal D}
\newcommand{\Glob}{\operatorname{Glob}}
\newcommand{\Lex}{\operatorname{Lex}}
\newcommand{\Set}{\mathbf{Set}}
\newcommand{\hhom}{\underline{\operatorname{Hom}}}
\title{Internal theories and intrinsic incompleteness}
\author{Florian Lengyel}
\date{29 October 2026}
\begin{document}
\maketitle
\begin{abstract}
An internal theory has an object of contradictions, obtained directly
as an internal hom-object. Ramesh's initiality construction equips
internal arithmetic universes with the introspection needed to apply
his L\"ob theorem. We deduce that the contradiction object of every
consistent internal arithmetic universe is an undecided proposition
in the initial arithmetic universe. Internal arithmetic functors
induce coherent comparison maps between these propositions and their
associated modal functors. This gives an intrinsic theory-and-extension
construction without choosing encodings into a universal object.
\end{abstract}

\section{The internal theory and its two modal functors}

An arithmetic universe is an exact, extensive category with
pullback-stable list objects. Arithmetic functors preserve this
structure. Write $\I$ for the initial arithmetic universe.
We use Ramesh's initiality construction and L\"ob theorem
\cite[Construction~6.11 and Theorem~4.19]{Ramesh}.

Let $\C$ be an arithmetic universe internal to $\I$.
Its external category $L_\C=\Glob(\C)$ has global objects
$c:1_\I\to\C_0$ and global arrows in $\C_1$ with the specified
endpoints. For global objects $c,d$, the object
$\hhom_\C(c,d)$ is the pullback of the source-target map
$\C_1\to\C_0\times\C_0$ along $(c,d):1_\I\to\C_0\times\C_0$.
In particular,
\begin{equation}\label{eq:global-hom}
 \I(1_\I,\hhom_\C(c,d))\cong L_\C(c,d).
\end{equation}
Internal composition induces composition on these objects.
The chosen internal arithmetic structure gives $L_\C$ its
arithmetic structure.

Initiality supplies an arithmetic functor
$S_\C:\I\to L_\C$, unique up to unique natural isomorphism.
The covariant functor
\[
 H_\C:L_\C\longrightarrow\I,
 \qquad H_\C(c)=\hhom_\C(1_\C,c)
\]
sends $u:c\to d$ to postcomposition by $u$.
It preserves finite limits, by the internal universal properties.
There are two covariant finite-limit-preserving endofunctors:
\begin{equation}\label{eq:two-boxes}
 B_\C=H_\C S_\C:\I\to\I,
 \qquad Q_\C=S_\C H_\C:L_\C\to L_\C.
\end{equation}
These act on different categories.

Ramesh's construction gives the natural transformation
$N_\C:\operatorname{id}_\I\Rightarrow B_\C$ required for introspection.
Its existence follows from the
initiality of the identity functor in $\Lex(\I,\I)$
\cite[Theorem~1.56 and Corollary~1.55]{Ramesh}.
Equivalently, consider the comma category $(\operatorname{id}_\I\mathbin\downarrow B_\C)$.
Arithmetic left comma-stability makes its projections arithmetic.
The initial arithmetic functor from $\I$ into this comma category
has both projections canonically isomorphic to $\operatorname{id}_\I$;
its comparison arrow is $N_\C$. The same initiality makes
this transformation unique.

The pair $(S_\C,N_\C)$ determines an indexed finite-limit-preserving
functor from $\I/{-}$ to the indexed category represented by $\C$
\cite[Theorem~2.5]{Ramesh}. Ramesh's L\"ob theorem therefore applies
to both functors in \eqref{eq:two-boxes}. In particular, an arrow
$B_\C A\to A$ gives a global element of $A$ in $\I$, and an
arrow $Q_\C c\to c$ gives a global element of $c$ in $L_\C$.
The fixed-point equations use the canonical identifications
$B_\C1_\I\cong1_\I$ and $Q_\C1_\C\cong1_\C$.

\section{The contradiction object}

A category with initial object $0$ and terminal object $1$ is
\emph{consistent} if it has no arrow $1\to0$.
A subterminal object $p$ is \emph{undecided} if neither
$1\to p$ nor $p\to0$ exists. A morphism $p\to0$ is a
refutation; no pseudocomplement object is required.

The category $\I$ is consistent: its arithmetic functor to $\Set$
would send an arrow $1_\I\to0_\I$ to a function from a singleton
to the empty set.

Define the contradiction object of $\C$ by
\begin{equation}\label{eq:contradiction}
 p_\C=\hhom_\C(1_\C,0_\C).
\end{equation}
This is subterminal in $\I$, since $0_\C$ is internally
subterminal. The arithmetic functor $S_\C$ preserves $0$ and $1$,
so $p_\C\cong B_\C0_\I$.

\begin{theorem}\label{thm:relative}
For every arithmetic universe $\C$ internal to $\I$, there is no
arrow $p_\C\to0_\I$.
Moreover, $p_\C$ is undecided in $\I$ if and only if
$L_\C=\Glob(\C)$ is consistent.
\end{theorem}
\begin{proof}
An arrow $p_\C\to0_\I$ would be an arrow
$B_\C0_\I\to0_\I$. L\"ob's theorem on $\I$ would give
$1_\I\to0_\I$, contradicting consistency of $\I$.
On the other hand, \eqref{eq:global-hom} gives
\[
 \I(1_\I,p_\C)\cong L_\C(1_\C,0_\C).
\]
Thus $p_\C$ has no global element precisely when $L_\C$ is
consistent. Together these prove the assertion.
\end{proof}

\begin{corollary}\label{cor:internal-consistency}
If $L_\C$ is consistent, there is no arrow
$Q_\C0_\C\to0_\C$ in $L_\C$.
\end{corollary}
\begin{proof}
Such an arrow would give $1_\C\to0_\C$ by L\"ob's theorem on
$L_\C$.
\end{proof}

Theorem~\ref{thm:relative} concerns $p_\C$ in $\I$;
Corollary~\ref{cor:internal-consistency} concerns
$Q_\C0_\C\cong S_\C(p_\C)$ in $L_\C$.
An arithmetic functor can send an object without global elements
to one having a global element. Consequently the corollary does
not assert that $S_\C(p_\C)$ is undecided in $L_\C$.

\section{Internal extensions}

Let $j:\C\to\D$ be an internal arithmetic functor.
It induces the arithmetic functor
$J=\Glob(j):L_\C\to L_\D$.
Initiality of $\I$ gives a canonical natural isomorphism
\[
 \lambda_j:JS_\C\xRightarrow{\ \cong\ }S_\D.
\]
Applying $j$ to internal arrows, and using its terminal-object
comparison and $\lambda_j$, defines
\[
 \beta_{j,A}: B_\C A\longrightarrow B_\D A
 \quad(A\in\I).
\]
In internal notation it sends $u:1_\C\to S_\C A$ to
$\lambda_{j,A}\,j(u)\,\epsilon_j^{-1}$, where
$\epsilon_j:j(1_\C)\xrightarrow{\cong}1_\D$.
Composition is written from right to left.

\begin{proposition}\label{prop:extensions}
The maps $\beta_{j,A}$ form a natural transformation
$\beta_j:B_\C\Rightarrow B_\D$ satisfying
$\beta_jN_\C=N_\D$.
Identity and composite internal arithmetic functors induce identity
and composite transformations, through the canonical comparisons.
Their components at $0_\I$ induce maps
$p_j:p_\C\to p_\D$.
If $L_\D$ is consistent, then $L_\C$ is consistent and both
$p_\C$ and $p_\D$ are undecided in $\I$.
\end{proposition}
\begin{proof}
For $f:A\to A'$ in $\I$, functoriality of $j$ and naturality
of $\lambda_j$ give the commutative square
\[
\begin{tikzcd}
B_\C A\arrow[r,"\beta_{j,A}"]\arrow[d,"B_\C f"']&
B_\D A\arrow[d,"B_\D f"]\\
B_\C A'\arrow[r,"\beta_{j,A'}"']&B_\D A'.
\end{tikzcd}
\]
Both $\beta_jN_\C$ and $N_\D$ are transformations
$\operatorname{id}_\I\Rightarrow B_\D$, so their equality follows from
the uniqueness established in Section~1.
The canonical comparison for a composite is the composite of
the canonical comparisons, by their uniqueness. The internal
functor laws then give the identity and composition assertions.
At $0_\I$, the map is the action of $j$ on
$\hhom_\C(1_\C,0_\C)$, with the initial- and terminal-object
comparisons. Finally, $J$ sends any arrow $1_\C\to0_\C$ to
an arrow $1_\D\to0_\D$. Apply Theorem~\ref{thm:relative}.
\end{proof}

These extensions can themselves be constructed by initiality.
Let $i:\mathbb T\to\mathbb T'$ be a finite-limit-preserving
map of finitely axiomatized lex theories. Write $\mathbb A$ for
the lex theory of arithmetic universes, and fix lex interpretations
$u:\mathbb A\to\mathbb T$ and $u':\mathbb A\to\mathbb T'$
with $iu=u'$.
Ramesh's internal-model theorem
\cite[Theorem~1.63 and Observation~1.66]{Ramesh} gives initial
models $G_{\mathbb T}:\mathbb T\to\I$ and
$G_{\mathbb T'}:\mathbb T'\to\I$.
Initiality gives the model homomorphism
\[
 G_{\mathbb T}\longrightarrow G_{\mathbb T'}i.
\]
On the underlying internal arithmetic universes it is an internal
arithmetic functor $j$. For successive extensions these maps compose
by uniqueness. Proposition~\ref{prop:extensions} therefore applies
to this explicitly specified class of theories and extensions.

\section{The initial example and the extension question}

Let $\C$ be the initial arithmetic universe internal to $\I$.
The external global-sections functor
$\Gamma_\I=\I(1_\I,-):\I\to\Set$ is arithmetic
\cite[Theorem~1.80]{Ramesh}. Indeed, it is initial among
finite-limit-preserving functors to $\Set$ by Yoneda; the unique
arithmetic functor to $\Set$ is also initial among these functors
by arithmetic left comma-stability. They are therefore naturally
isomorphic. Preservation of the internal initial model then gives
$\Glob(\C)\simeq\I$ and makes $S_\C$ an equivalence
\cite[Observation~6.12]{Ramesh}.
Thus $\Glob(\C)$ is consistent, and Theorem~\ref{thm:relative}
gives a subterminal $p_\C$ isomorphic to neither $0_\I$ nor
$1_\I$, with neither a proof nor a refutation in $\I$.

The external arithmeticity of $\Gamma_\I$ does not internalize.
For this very example, the internal global-sections functor sends
$0_\C$ to $p_\C$, which is not $0_\I$.
This distinction accounts for the simultaneous external consistency
and internal failure to refute $p_\C$.

The remaining essential-incompleteness question is to prove, for the
intended consistent internal extensions $\D$, that their own
category $L_\D$ has an undecided proposition. The results above
construct undecided propositions in the fixed base $\I$ and prove
the second-incompleteness assertion in each consistent $L_\D$.
Passing from these statements to undecidability in $L_\D$ requires
a further argument; the equivalence $\Glob(\C)\simeq\I$ for the
internal initial example supplies no such equivalence for general
extensions.

\begin{thebibliography}{9}
\bibitem{Ramesh} S.~Ramesh,
\emph{Introspective Theories and Geminal Categories},
Ph.D. dissertation, University of California, Berkeley, 2023.
\url{https://escholarship.org/uc/item/3mn0c475}.
\end{thebibliography}
\end{document}
